\newpage
\phantomsection
\label{prf:rational-addition-well-defined}

\begin{remark*}[Return]
\hyperref[lem:rational-addition-well-defined]{Return to Lemma}
\end{remark*}

\begin{lemma*}[Addition Is Well-Defined on Rational Classes]
Let $a,b,c,d,e,f,g,h\in\mathbb{Z}$ with $bdfh\neq 0$. If
\[
(a,b)\sim(e,f)
\qquad\text{and}\qquad
(c,d)\sim(g,h),
\]
then
\[
(ad+bc,bd)\sim(eh+fg,fh).
\]
\end{lemma*}

\begin{proof}
\textbf{Professional Standard Proof.}~\\
Assume
\[
(a,b)\sim(e,f)
\qquad\text{and}\qquad
(c,d)\sim(g,h).
\]
By the definition of $\sim$,
\[
af=be
\qquad\text{and}\qquad
ch=dg.
\]
To prove
\[
(ad+bc,bd)\sim(eh+fg,fh),
\]
it is enough to prove
\[
(ad+bc)fh=bd(eh+fg).
\]

From $af=be$, multiplying both sides by $dh$ gives
\[
adfh=bdeh.
\]
From $ch=dg$, multiplying both sides by $bf$ gives
\[
bcfh=bdfg.
\]
Adding these two equalities gives
\[
adfh+bcfh=bdeh+bdfg.
\]
Factoring both sides gives
\[
(ad+bc)fh=bd(eh+fg).
\]
Therefore
\[
(ad+bc,bd)\sim(eh+fg,fh).
\]
Hence addition is well-defined on rational classes.
\end{proof}

\begin{proof}
\textbf{Detailed Learning Proof.}~\\

\textbf{Step 1.} Translate the hypotheses into cross-product equalities.

Assume
\[
(a,b)\sim(e,f)
\qquad\text{and}\qquad
(c,d)\sim(g,h).
\]
By the definition of the fraction-pair relation, these assumptions mean
\[
af=be
\qquad\text{and}\qquad
ch=dg.
\]

\textbf{Step 2.} Translate the desired conclusion into one equality.

The desired conclusion is
\[
(ad+bc,bd)\sim(eh+fg,fh).
\]
By the definition of $\sim$, this is equivalent to the cross-product equality
\[
(ad+bc)fh=bd(eh+fg).
\]

\textbf{Step 3.} Produce the first summand equality.

From
\[
af=be,
\]
multiplying both sides by $dh$ gives
\[
adfh=bdeh.
\]
This matches the first summand on each side of the target equality.

\textbf{Step 4.} Produce the second summand equality.

From
\[
ch=dg,
\]
multiplying both sides by $bf$ gives
\[
bcfh=bdfg.
\]
This matches the second summand on each side of the target equality.

\textbf{Step 5.} Add and factor.

Adding the two equalities
\[
adfh=bdeh
\qquad\text{and}\qquad
bcfh=bdfg
\]
gives
\[
adfh+bcfh=bdeh+bdfg.
\]
Factoring the left-hand side and the right-hand side gives
\[
(ad+bc)fh=bd(eh+fg).
\]
Therefore, by the definition of $\sim$,
\[
(ad+bc,bd)\sim(eh+fg,fh).
\]
\end{proof}

\begin{remark*}[Proof structure]
The proof is direct. The hypotheses are first converted into the two
cross-product equalities defining equivalence of fraction pairs. Each
equality is then multiplied by the remaining denominators needed to
match the common target denominator. Adding the resulting equalities
and factoring gives the cross-product equality required for the sums.
\end{remark*}

\begin{remark*}[Dependencies]~\\
\begin{itemize}
  \item \hyperref[def:fraction-pair-equality]{Definition~\ref*{def:fraction-pair-equality}}
  \item Basic arithmetic of the integers
  \item Distributivity of integer multiplication over addition
\end{itemize}
\end{remark*}
